\section{Verification purpose}

This case verifies temporal-resolution sensitivity using a fixed physical
problem and four level-set timesteps. Every curve identified as ELMFIRE output
is regenerated from case-specific simulation products.

This case verifies temporal-resolution sensitivity in a firebrand-driven WUI calculation. Spatial resolution, geometry, wind, heat-release history, firebrand physics, random seed, and output cadence are fixed; only the level-set timestep changes over a wide Courant--Friedrichs--Lewy (CFL) range.

The test determines the range over which ignition chronology and structure-level spread remain timestep independent and makes loss of resolution explicit at very large steps. Success is evaluated against the finest-step baseline using load, ignition-time, and spread metrics.

\section{Mathematical formulation and numerical implementation}

\subsection{Fixed physical problem}

The community is a \(220\times100\)-m two-dimensional lattice of
\SI{10}{m} square structures separated by \SI{10}{m} inert firebreaks in both
directions. Every structure in the first upwind column is burning at \(t=0\), and wind is uniform at
\SI{40}{mph} from west to east. Non-structure cells use nonburnable fuel model
93, so vegetative surface-fire spread is absent. Thus a
new source can appear only after deposited firebrands complete the configured
physical ignition process.

For a structure ignited at \(t_i\), the design heat-release rate (HRR) per unit area is
\[
q''(\tau)=
\begin{cases}
q''_{\max}\tau/300,&0<\tau<300,\\
q''_{\max},&300\le\tau\le3900,\\
q''_{\max}(4200-\tau)/300,&3900<\tau<4200,\\
0,&\text{otherwise},
\end{cases}
\quad \tau=t-t_i,\qquad q''_{\max}=\SI{400}{kW.m^{-2}}.
\]
For footprint area \(A_s=\SI{100}{m^2}\),
\(\mathrm{HRR}_i=q''A_s\). Structural firebrands are emitted using the per-power rate
\[
\dot N_i(t)=GR'\times{\mathrm{HRR}_i(t)},\qquad
GR'=\SI{10}\,\rm{pcs\, MW^{-1}\,s^{-1}}.
\]

\subsection{Transport, deposition, and analytical load}

The downwind landing distance \(X\) is represented by the lognormal density
\[
 f(x\mid\mu,\sigma)=
 \frac{\exp[-(\ln x-\mu)^2/(2\sigma^2)]}
      {x\sigma\sqrt{2\pi}},\qquad x>0.
\]
For source structure \(i\), the heat-release-dependent transport parameter is
\[
 B_i^*=\frac{U}{\sqrt{gL}}
 \left(\frac{\rho_f}{\rho_a}\right)^{-3/4}
 \left(\frac{d_f}{L}\right)^{-3/4}
 \left[
 \frac{\mathrm{HRR}_i}
 {\rho_a c_p T_a\sqrt{g}\,L^{5/2}}
 \right]^{1/2}.
\]
The dimensional mean and standard deviation of landing distance are
\[
 \bar x_i=0.47(B_i^*)^{2/3}L,\qquad
 s_{x,i}=0.88(B_i^*)^{1/3}L,
\]
and the log-space parameters supplied to \(f\) are
\[
 \sigma_i=\sqrt{\ln\!\left[1+(s_{x,i}/\bar x_i)^2\right]},
 \qquad
 \mu_i=\ln\!\left(
 \frac{\bar x_i}{\sqrt{1+(s_{x,i}/\bar x_i)^2}}
 \right).
\]
Here \(L=\SI{10}{m}\), \(d_f=\SI{5e-3}{m}\),
\(\rho_f=\SI{100}{kg.m^{-3}}\), \(\rho_a=\SI{1.1}{kg.m^{-3}}\),
\(c_p=\SI{1}{kJ.kg^{-1}.K^{-1}}\), \(T_a=\SI{300}{K}\),
\(g=\SI{9.81}{m.s^{-2}}\), and
\(U=0.447(40)/0.87=\SI{20.55}{m.s^{-1}}\). This parameterization and its
configured constants follow the documented structural-firebrand
model~\cite{Qin2025}; the equations are reproduced here so the reference can
be evaluated without consulting that source.

The independently calculated deposited count on target structure \(k\) is
\[
N_k=\sum_{i<k}\int_0^T GR'\,\mathrm{HRR}_i(t)
\left[\int_{x_k}^{x_{k+1}}
 f\!\left(x-x_i\mid\mu_i(t),\sigma_i(t)\right)\,dx\right]dt .
\]
The inner integral is evaluated using the lognormal cumulative distribution function (CDF) and the time integral
uses an independent \SI{5}{s} midpoint quadrature. ELMFIRE's final deposited
count is divided by the physical \SI{100}{m^2} structure footprint, giving
\(N_k/A_s\) in \si{pcs\,m^{-2}}.

The Eulerian firebrand tracker advances with its own wind-limited substeps.
Changing the level-set step therefore must not change the spatial resolution of
the flight trajectory. This is the core algorithmic separation being tested.
In the generated namelists, prescribed timestep and
maximum timestep impose the four values above. The input
target surface-spread Courant number remains the ordinary surface-spread limiter and does not
redefine the four wind-based temporal simulation conditions.  The global vegetative surface-spread setting switch remains true to suppress vegetative surface-fire propagation.  The current source independently refreshes initialized model-14 building HRR, and the nonburnable fuel-93 background provides an additional isolation of building-only spread.
The fractional timesteps are intrinsic to this temporal-resolution sweep.  For
each member, preprocessing sets the generated stop to
\(\lceil t_{\mathrm{requested}}/\Delta t\rceil\Delta t\), and records the
requested stop, aligned stop, and integer step count in the simulation specification.

\subsection{Physical ignition}

For the configured structural ignition model, the critical deposited load is
\(N_{\mathrm{crit}}=\SI{2928.9}{pcs\,m^{-2}}\), the
smoldering-to-flaming timescale is \SI{42.1}{s}, and the subsequent structure
development delay is \SI{300}{s}~\cite{Qin2025}. Because these are much longer
than the CFL=0.5, 1, and 10 level-set steps, their ignition chronologies should
remain nearly unchanged. CFL=200 deliberately makes the time step comparable
to or larger than the short ignition time scale and exposes the onset of
temporal sensitivity.

\section{Derivation of expected behavior}
When firebrand flight and ignition time scales are much longer than the level-set step, moderate timestep changes should preserve deposited-load histories and structure ignition times. Extremely large steps should show increasing temporal error without changing the fixed physical problem.

% BEGIN EXPLICIT EXPECTED-RESULT DERIVATION
At peak HRR, one \(10\times10~\mathrm{m}\) structure produces
\[
 HRR_{\max}=400(100)=40~\mathrm{MW},\qquad
 \dot N_{\max}=10(40)=400~\mathrm{pcs\,s^{-1}},
\]
and the critical total load is
\(2928.9(100)=292890\) particles.  The transport reference uses the adjusted
air speed \(0.447(40)/0.87=20.55~\mathrm{m\,s^{-1}}\), whereas preprocessing
constructs the declared level-set timesteps from the 20-ft input speed
\(U_{20}=40(0.44704)=17.8816~\mathrm{m\,s^{-1}}\).  Thus
\[
\begin{array}{c|rrrr}
CFL_w&0.5&1&10&200\\\hline
\Delta t=CFL_w(10)/17.8816~(\mathrm{s})
&0.279617&0.559234&5.592341&111.846815
\end{array}
\]
which reproduces the simulation condition names to their displayed precision.

Relative to the \(\SI{42.1}{s}\) small-fuel timescale, these steps are
\(0.00664,\ 0.0133,\ 0.1328,\) and \(2.657\) times as large.  The first
three therefore resolve that timescale, while CFL 200 spans more than twice
it and is expected to show temporal degradation.  The analytical convolution
is unchanged among simulation conditions, so its ideal load error is zero.  The explicit
numerical substitutions are: baseline mean load error at most 0.01; CFL 1
and 10 maximum load and mean ignition-time errors at most 0.001; their
history NRMSE at most 0.05; and CFL 200 maximum load error at most 0.10.
% END EXPLICIT EXPECTED-RESULT DERIVATION

\section{Verification metrics and acceptance criteria}

Prior calculations for this fixed physical problem motivate the following
predeclared expectations~\cite{Qin2025}:
\begin{itemize}
\item the maximum analytical-load error remains below 10\% even at CFL=200;
\item for CFL no greater than 10, analytical-load errors are generally below
0.1\%;
\item ignition-time errors relative to CFL=0.5 remain below 0.1\% for CFL no
greater than 10;
\item accumulation histories at the structure spanning
\(140\le x\le\SI{150}{m}\) closely align; and
\item first-deposition timing remains close through the approximately
\SI{80}{m} peak-HRR spotting range.
\end{itemize}

The case converts those expectations into explicit acceptance checks. The
CFL=0.5 mean error against the analytical deposited-load integral defined
above must not exceed 1\%; the CFL=1 and 10 maximum errors against the same
integral must not exceed 0.1\%; and the CFL=200 maximum must remain below
10\%. CFL=1 and 10 mean full-ignition-time errors relative to CFL=0.5 must not
exceed 0.1\%. Accumulation-history normalized RMSE at 140--150~m must not
exceed 5\% for CFL=1 and 10. Every simulation condition must be complete, and every
acceptance component must pass.

The postprocessor also estimates first-deposition time from the first transient
output interval with nonzero accumulation. Because its resolution is the
\SI{500}{s} dump interval rather than the firebrand trajectory substep, this is
shown as a diagnostic corresponding to the first-deposition diagnostic, but is not used as a
pass/fail metric.

\subsection{Justification of metric quantification}

The 0.1\% limits for CFL 1 and 10 and the 10\% limit for CFL 200 quantify the
reported temporal-resolution behavior for this fixed-grid experiment rather
than being selected from the present outputs~\cite{Qin2025}.  They deliberately
ask for near invariance over the moderate-step range while treating the
extreme-step case as acceptable only if its error remains bounded rather than
converged.  The baseline mean-load limit is 1\% because it compares a
cell-integrated analytical deposition relation with a thresholded numerical
field; it establishes that the reference member itself is accurate enough to
anchor the relative ignition comparisons.

The 5\% history NRMSE limit is wider than the final-load limits because
transient rasters are sampled every \SI{500}{s}; small timing offsets at a
rising accumulation front are amplified pointwise even when final totals
agree.  Normalizing by the baseline history range makes 5\% a direct bound on
the fraction of the resolved accumulation excursion.  First-deposition time is
kept diagnostic because its \SI{500}{s} observation interval is too coarse to
support a defensible tighter decision.  All four simulation conditions are required so the
case demonstrates both moderate-step invariance and the predicted extreme-CFL
degradation.

\subsection{Predeclared acceptance criteria}

Table~\ref{tab:acceptance-contract} summarizes the metrics fixed before
inspection of the calculated results. Numerical values and component statuses
are reported separately in Section~7.

\begin{table}[H]
\centering
\footnotesize
\caption{Predeclared verification metrics and acceptance criteria.}
\label{tab:acceptance-contract}
\begin{tabular}{@{}p{0.23\linewidth}p{0.47\linewidth}p{0.21\linewidth}@{}}
\toprule
Metric & Output, region, and aggregation & Acceptance criterion\\
\midrule
Output completeness & Current-input identity record outputs for baseline and CFL 1, 10, and 200 simulation conditions. & 4 of 4 simulation conditions \\
Baseline load mean error & Mean structure-load relative error against the analytical load reference. & \(\le0.01\) \\
CFL 1 load maximum error & Maximum finite structure-load relative error against its analytical reference. & \(\le0.001\) \\
CFL 10 load maximum error & Same maximum error at CFL 10. & \(\le0.001\) \\
CFL 200 load maximum error & Same maximum error at CFL 200. & \(\le0.10\) \\
CFL 1 ignition error & Mean matched full-ignition-time relative error against CFL 0.5. & \(\le0.001\) \\
CFL 10 ignition error & Mean matched full-ignition-time relative error against CFL 0.5. & \(\le0.001\) \\
CFL 1 history NRMSE & Normalized RMSE of accumulation history at the 140--150-m structure against baseline. & \(\le0.05\) \\
CFL 10 history NRMSE & Same normalized history error at CFL 10. & \(\le0.05\) \\
Overall decision & Conjunction of completeness and all eight quantitative checks. & PASS only if all pass \\
\bottomrule
\end{tabular}
\end{table}

\section{Simulation configuration and predicted outcome}

Figure~\ref{fig:input-configuration} records the prepared spatial inputs for the wind Courant number 0.5 and timestep 0.28 s condition.
These panels show the prepared spatial fields, remove the
two-cell numerical buffer, and retain map coordinates in metres.  The left panel shows
the most informative fuel or structure layout available for the case; the right panel
shows the cells selected by the non-positive initial level-set field, making a localized
ignition or firebrand-emission source explicit.

\begin{figure}[H]
\centering
\EvidenceFigure[width=0.98\textwidth,height=0.42\textheight,keepaspectratio]{../figures/input_configuration.pdf}
\caption{Whole-domain prepared input configuration for the representative simulation condition.  The
physical domain is shown without the numerical halo; coordinates, configured wind values,
fuel or structure layout, and initial ignition/source geometry are identified in the panels.}
\label{fig:input-configuration}
\end{figure}

\begin{table}[H]
\centering\small
\caption{Generated temporal-resolution simulation conditions. The time step is calculated
from \(\Delta t=\mathrm{CFL}\,\Delta x/U\), not rounded to the display value.}
\begin{tabular}{lrrr}
\toprule
Simulation condition & Wind-based CFL & \(\Delta t\) (s) & \(t_{\max}\) (s)\\
\midrule
\VariantRows
\bottomrule
\end{tabular}
\end{table}

All simulation conditions use \(\Delta x=\Delta y=\SI{10}{m}\), alternating
\SI{10}{m} structures and firebreaks, a \SI{40}{mph} wind, the same model-14
design fire, per-MW firebrand generation, empirical structural transport,
Eulerian accumulation, physical ignition, no ember consumption, no crosswind
distribution, and no vegetative surface-fire pathway. The predicted result is near-overlap for
CFL=0.5, 1, and 10, with visible but bounded load sensitivity at CFL=200.

\subsection{Preprocessing, execution, and postprocessing}

Each prescribed condition uses co-registered fields in a Universal Transverse
Mercator coordinate system and the stated fuel properties. The initial level-set
field \(\phi\) is negative within every structure footprint in the first upwind column.
The exact timestep and input identities are recorded for each condition. Unique
structure identities preserve the two-dimensional lattice; the profiles compare
a representative cross-stream row.

A completed result is reused only when its input identity matches the current
prescribed configuration. Each simulation condition is evaluated independently.

result analysis procedure aggregates final ember counts and ignition times
by physical structure ID, reconstructs cumulative histories from transient
ember rasters and recorded output times metadata, evaluates the analytical deposited-load relation,
calculates every declared metric, and produces report-ready vector figures.
Stale pre-modification outputs are explicitly excluded.

\section{Actual simulation results}

Figure~\ref{fig:domain-result} shows the representative time-of-arrival field from the wind Courant number 0.5 and timestep 0.28 s condition.
The displayed field is taken from the simulated results, masks missing values, and excludes
the two-cell numerical buffer.  The finest-timestep reference field anchors the profile and error comparisons and exposes domain-wide behavior that a single-row profile can hide.

\begin{figure}[H]
\centering
\EvidenceFigure[width=0.96\textwidth,height=0.50\textheight,keepaspectratio]{../figures/domain_result.pdf}
\caption{Whole-domain time-of-arrival field from the representative completed ELMFIRE run.  This
domain-scale result supplements, rather than replaces, the higher-level verification
profiles, convergence plots, ensemble statistics, and calculated acceptance metrics below.}
\label{fig:domain-result}
\end{figure}

\CompletedVariantCount{} of \RequestedVariantCount{} simulation conditions currently have
complete outputs matching the generated inputs. The verification decision is
\textbf{\OverallStatus}.

Figure~\ref{fig:temporal-four-panel} presents four complementary
temporal-resolution diagnostics. Panel (a) compares deposited load with the
analytical integral defined in Section~2; panels (b) and (d) use CFL=0.5 as
the numerical baseline; and panel (c) shows the accumulation history at the
specified downwind structure.

\begin{figure}[H]
\centering
\EvidenceFigure[width=\linewidth]{temporal_resolution_profiles.pdf}
\caption{Temporal-resolution diagnostics for the prescribed temporal-resolution sweep:
(a) accumulated-load relative error; (b) dump-cadence first-deposition-time
relative error; (c) accumulation history at 140--150~m; and (d) full-ignition
time relative error.}
\label{fig:temporal-four-panel}
\end{figure}

Figure~\ref{fig:mean-ignition-error} condenses the ignition comparison into
mean relative error versus CFL and displays the acceptance threshold directly.

\begin{figure}[H]
\centering
\EvidenceFigure[width=0.72\linewidth]{temporal_resolution_error.pdf}
\caption{Mean full-ignition-time relative error referenced to CFL=0.5. The
dashed line is the 0.1\% acceptance level used for CFL=1 and 10.}
\label{fig:mean-ignition-error}
\end{figure}

\section{Calculated metrics and verification decision}

\begin{table}[H]
\centering\small
\caption{Calculated verification metrics and component decisions.}
\begin{tabular}{@{}p{0.37\linewidth}p{0.18\linewidth}p{0.17\linewidth}p{0.14\linewidth}@{}}
\toprule
Metric & Acceptance limit & Calculated & Status\\
\midrule
\MetricRows
\bottomrule
\end{tabular}
\end{table}

The current decision is \textbf{\OverallStatus}. A pass requires all four
input-matched runs and all tabled acceptance checks. A not-run state is
not evidence of agreement.

This is a deterministic code-verification test of the implemented idealized WUI
model, not field validation. The analytical load uses simulated source ignition
times so that it isolates generation and deposition from ignition chronology.
First-deposition timing is limited by transient output cadence because the
current executable does not write the internal firebrand arrival time array as a
final raster; that diagnostic is therefore intentionally excluded from the
decision specification.

\section{Technical interpretation}
With the declared building-only switch restored, all four current-source simulation conditions produce nonzero structural heat release, firebrand transport, input identity record-matched final rasters, and a complete metric record.  This invalidates the earlier zero-HRR and missing-output interpretations, both of which arose while diagnosing the stale executable or the temporary altered switch.

The temporal history is comparatively stable at moderate step sizes: the CFL 1 and CFL 10 accumulation-history NRMSE values are 0.008466 and 0.008443, both below 0.05.  The stricter pointwise/reference components do not pass.  The analytical-load errors are 0.54345 for CFL 0.5, 0.82701 for CFL 1, 0.82690 for CFL 10, and 0.83927 for CFL 200, all above their respective limits.  The CFL 1 and CFL 10 mean ignition-time relative errors are 0.007295 and 0.01326, compared with the 0.001 limit.  The overall result is therefore \emph{FAIL}, not \emph{NOT RUN} or \emph{NOT EVALUABLE}.

Because every input identity record, final output, and integer timestep/stop relation is valid, these discrepancies are not caused by an admissible case-configuration error.  They identify disagreement between the implemented temporal/structural spotting calculation and the predeclared analytical-load and ignition-time requirements.

\section{Scope and limitations}
The test varies only temporal resolution in a fixed academic WUI configuration. It verifies timestep sensitivity for the selected transport and ignition scales and does not establish spatial convergence or physical validation.

\begin{thebibliography}{9}
\bibitem{Qin2025}
Y. Qin, \emph{A Physics-Based Eulerian Framework for Modeling Firebrand
Showering in Regional-Scale Wildland and Wildland--Urban Interface Fire
Simulations}, Ph.D. dissertation, University of Maryland, College Park, 2025.
\end{thebibliography}
